\chapter{Ковариантная фиксация калибровки и отрицательная мода вихря}

\section{Диагноз мягкого сектора}

Предыдущая глава обнаружила восемь уровней, смягчающихся при сгущении
сетки, но не смогла отождествить их с переносами. Разложение собственных
векторов показывает, что около половины их нормы лежит в заряженной
амплитуде, около половины --- в связности, а вклад нейтрального поля
меньше $10^{-5}$. Это указывает прежде всего на проверку знака фоновой
калибровки.

При соглашении
\begin{equation}
 D_i\phi=\partial_i\phi-A_iJ\phi
\end{equation}
эллиптическое условие имеет вид
\begin{equation}
 \mathcal G_-(\eta,\alpha)
 =\partial_i\alpha_i-\frac AG(J\phi_0)\mathbin{\cdot}\eta.
 \label{eq:v6-covariant-minus-gauge}
\end{equation}
На калибровочной орбите соответствующий оператор Фаддеева--Попова имеет
положительную эллиптическую часть
\begin{equation}
 -\Delta+\frac AG|\phi_0|^2.
\end{equation}
Знак плюс в предыдущем условии давал вместо неё оператор типа
$\Delta+(A/G)|\phi_0|^2$ и допускал искусственное приближение
калибровочных уровней к нулю.

К физической второй вариации добавлен неотрицательный квадрат
\begin{equation}
 E_{\mathrm{fix}}=\frac G2\int_{\mathbb R^2}\mathcal G_-^2\,d^2x.
 \label{eq:v6-positive-gauge-fixing-square}
\end{equation}
Поэтому отрицательное значение оператора с этим добавленным квадратом не
может быть создано фиксацией калибровки: исходная физическая квадратичная
форма на том же направлении ещё меньше.

Такое требование согласуется с анализом нулевых мод БПС-вихрей в
\path{arXiv:1602.09084} и с самосопряжённой ортогональной калибровкой
\path{arXiv:2608.10608}.

\section{Изолированная отрицательная мода}

Исправленный оператор пересчитан в квадрате радиуса $10$. Нижнее
собственное значение на сетках $24,32,40,48$ равно соответственно
\begin{equation}
 0.01510310,\quad -0.10337578,\quad
 -0.20169832,\quad -0.18904743.
 \label{eq:v6-negative-mode-convergence}
\end{equation}
Сдвиг между двумя последними сетками равен $0.01265089$. Следующий уровень
на сетке $48$ положителен и равен $3.4097\cdot10^{-4}$, так что
отрицательная мода изолирована от мягкого калибровочно-переносного
подпространства.

Проверка размера области при сетке $40$ сохраняет знак:
\begin{equation}
 \lambda_{\min}(R=8)=-0.18007196,qquad
 \lambda_{\min}(R=12)=-0.13534038.
\end{equation}
На самой тонкой основной сетке отрицательный вектор распределён между
заряженной амплитудой и связностью:
\begin{equation}
 w_\phi=0.521256,qquad
 w_A=0.476002,qquad
 w_b=0.002742.
\end{equation}
Следовательно, это не чистая калибровочная мода и не колебание
нейтрального профиля, а совместная деформация потока и заряженного ядра.

Для прямой нити продольный импульс даёт
\begin{equation}
 \lambda(k_z)=\lambda(0)+k_z^2.
\end{equation}
При значении сетки $48$ неустойчива полоса
\begin{equation}
 |k_z|<\sqrt{-\lambda(0)}=0.43480\ldots.
\end{equation}
Это длинноволновая неустойчивость прямой нити, а не ультрафиолетовый
дефект дискретизации.

\section{Физический смысл и граница}

Канонический прямой $Z_6$-вихрь является седлом уже внутри эффективной
подсистемы. Тем самым он не может служить устойчивой бесконечной
материальной нитью. Однако отрицательная мода не закрывает рождение
материи: пока неизвестно, к чему она нелинейно ведёт --- к распаду,
расщеплению ядра, скручиванию, замыканию нити или другому дефекту.

\begin{nogocriterion}
Запрещено называть отрицательную моду хопфовым замыканием или рождением
частицы до нелинейного продолжения вдоль собственного вектора и
вычисления конечной энергии нового стационарного состояния. Если энергия
монотонно убывает к распаду, прямой вихрь закрывается без материального
остатка.
\end{nogocriterion}

\begin{protocol}
\begin{enumerate}
 \item знак фоновой калибровки: \textbf{исправлен};
 \item остаточный калибровочный мягкий кластер: \textbf{диагностирован};
 \item изолированная отрицательная мода: \textbf{найдена};
 \item минимальное значение: \textbf{$-0.18904743\ldots$};
 \item прямая вихревая нить: \textbf{линейно неустойчива};
 \item нелинейный исход: \textbf{не определён};
 \item следующий гейт: \textbf{нелинейное насыщение отрицательной моды}.
\end{enumerate}
\end{protocol}

Машинный сертификат сохранён в
\path{s2t_v6_bosonic_defect_corrected_vortex_covariant_zero_mode_resolution_gate_results.json}.